\chapter{Bibliographie}

\section*{Sentiments et communautés patrimoniales}

\subsection*{Monographies}
Bermès Emmanuelle, \guillemetleft De l'écran à l'émotion : Quand le numérique devient patrimoine \guillemetright, éditions de \textit{l'École nationale des chartes - PSL}, 242 pages, Paris, 2024 \\

De Kosnik Abigail, \guillemetleft Rogue Archives, Digital Cultural Memory and Media Fandom, éditions \textit{The MIT Press}, 430 pages, Cambridge, 2016 \\

Fabre Daniel, \og Émotions patrimoniales \fg, collection Ethnologie de la France,\textit{ éditions de la Maison des sciences de l’homme}, 2013. \\

Gray Jonathan, Sandvoss Cornel, Harrington C. Lee, \og Fandom: identities and communities in a mediated world \fg, éditions \textit{New York University Press}, 2017.\\

\subsection*{Articles de revues}
Blanchard Jean-François, \og Chiara Bortolotto, Le patrimoine culturel immatériel. Enjeux d’une nouvelle catégorie \fg, éditions \textit{ENS Editions}, 

Lessig Lawrence, \guillemetleft Industrie de la culture, pirates et \guillemetleft culture libre\guillemetright \guillemetright, dans \textit{Critique}, volume 73, n°6, p510-518, 2008\\

Paloque-Berges Camille, \og Vers des lieux de mémoire réticulaires ? \fg, dans \textit{Reset, Recherches en sciences sociales sur l'internet}, n°6, 2016. 

Paloque-Berges Camille et Schafer Valérie, \og Quand la communication devient patrimoine \fg, dans la Revue Hermès, \textit{éditions du CNRS}, p. 255-261, 2015. \\

\section*{Piratage et culture pirate}

\subsection*{ Monographies}

Baumgärtel Tilman, \guillemetleft A Reader on International Media Piracy : Pirate Essays \guillemetright, collection \textit{MediaMatters Series}, éditions \textit{Amsterdam University Press}, 2015\\

Paul Eve Martin, \og Warez, The Infrastructure and Aesthetics of Piracy \fg, éditions punctumbooks, 445 pages, 2021.\\

Scorpecci Danny et Stryszowski Piotr , \og Piracy of Digital Content \fg, éditions \textit{OECD Publishing}, 2009.

\subsection*{Articles de revues }

Bohannon John, \og Who's downloading pirated papers ? Everyone \fg, dans \textit{Science}, vol 352, p. 508-512, 2016. \\

Galluzzo Anthony, \og Longévité et résilience de l’accès illégal aux contenus culturels. Interroger les persistances et les sorties de carrières pirates \fg, dans \textit{Annales des Mines \- Gérer \& comprendre}, volume 145 n°3, p. 27-45, 2021. \\

Latrive Florent, \guillemetleft Du bon usage de la piraterie. Culture libre, science ouverte~\guillemetright, dans \textit{Multitude} volume 18, n° 4, p. 197-202, 2004 \\

Liang Lawrence, \guillemetleft Piratage, créativité et infrastructure : repenser l'accès à la culture \guillemetright, dans \textit{Tracés, revue de sciences humaines}, p183-202, 2014\\

Mattelart Tristan, \guillemetleft Piratage : apport et limites d'une infrastructure d'accès à la culture \guillemetright, dans \textit{Tracés, revue de sciences humaines}, n°26, p175-182, 2014\\

Tatchim Nicanor, \og Piratage, économie informelle et industries culturelles au Cameroun : approches socio-économique, communicationnelle et politique \fg, dans \textit{Études de communication, Médiations et usages sociaux des savoirs et de  l’information. Regards croisés France-Brésil}, n°57, 2021

Taylor Diana, \og Save As ... Knowledge and Transmission in the Age of Digital Technologies \fg, dans Imagining America n°7, 2010. 

Todd Darren, \og Pirate nation: how digital piracy is transforming business, society and culture \fg, éditions \textit{Kogan Page}, 2011. 

\section*{Archives et bibliothèques sur le web}

\subsection*{Monographies}

Brügger Niels, \og The archived web : doing history in the digital age \fg, éditions \textit{The MIT Press}, 2018. 

\subsection*{Articles de revues}

Magis Crisophe et Granjon Fabien, \og Numérique et libération de la produciton scientifique \fg, dans \textit{Variations. Revue internationale de théorie critique}, n°19, 2016. \\

Paul Eve Martin, \guillemetleft Lessons from the library : Extreme Minimalist Scaling at Pirate Ebook Platforms\guillemetright, dans \textit{Digital Humanities Quarterly}, dans la revue \textit{Providence}, volume 16, n° 2, 2022\\

\section*{Propriété intellectuelle et artistique}

\subsection*{Monographies}

National Research Council, \og The digital dilemma: intellectual property in the information age\og, \textit{National Academy Press}, 2010 \\

Sell Susan K,\og Private powern public law : the globalization of intellectual property rights \fg, collection Cambridge studies in international relations, éditions \textit{Cambridge University Press}, 2009

Goldstein Paul, \og, International Copyright : principles, law and practices \fg, \textit{Oxford University Press}, 2001. 